\documentclass[10pt,a4paper]{amsart}
\usepackage[margin=19mm]{geometry}
\usepackage{amsmath,amssymb,mathtools}
\usepackage[scr=rsfs,cal=euler]{mathalfa}
\usepackage{xfrac}
\usepackage[unicode,hidelinks,bookmarks=false]{hyperref}
\newcommand{\ie}{i.e.\ }
\newcommand{\cf}{cf.\ }
\newcommand{\resp}{respectively\ }
\newcommand{\Subsection}{\subsection}
\providecommand{\nobreakdash}{\nobreak\hbox{-}}
\setlength{\emergencystretch}{2em}
\multlinegap0pt
\usepackage{bm}
\usepackage{bbm}
\usepackage{dsfont}
\usepackage{xspace}
\usepackage{cancel}
\usepackage{mathtools}
\usepackage{amsthm}
\makeatletter
\@ifundefined{Definition}{\newtheorem{Definition}{Definition}}{}
\@ifundefined{Lemma}{\newtheorem{Lemma}[Definition]{Lemma}}{}
\@ifundefined{Proposition}{\newtheorem{Proposition}[Definition]{Proposition}}{}
\makeatother
\title[Moser's twist theorem revisited]{Moser's twist theorem revisited}
\author{Yi Liu, Lin Wang}
\date{Source: arXiv:2502.10969v3 (CC BY 4.0)}
\begin{document}
\maketitle

\noindent\emph{Source and licence.} Moser's twist theorem revisited. Version arXiv:2502.10969v3 (CC BY 4.0), distributed under the Creative Commons Attribution 4.0 International licence, as recorded in the arXiv licence metadata for this exact version and shown on the abstract page https://arxiv.org/abs/2502.10969v3. Full rights notice: licenses/arxiv-2502.10969-CC-BY-4.0.txt.\par\smallskip

\noindent\textit{Editorial scope.} Counted unit: the Lemma whose source label is \texttt{karrr1} in the article (source0.tex lines 954--960, source label \texttt{karrr1}), delivered together with the article's own complete original proof of it (source0.tex lines 961--1063, one \texttt{proof} environment of 103 lines). No mathematics is written, repaired or re-derived: statement and proof are the article's own text line for line. Required context retained uncounted: Z12 (lines 404--406); ko11 (lines 604--606); f2x (lines 670--672); K22 (lines 830--833); keyy1 (lines 836--850); r222 (lines 931--933); arthh (lines 937--939); KES (lines 949--951). Every definition, display and label of those lines is retained unchanged. Omitted from this package and not redistributed: 1--403, 407--603, 607--669, 673--829, 834--835, 851--930, 934--936, 940--948, 952--953, 1064--1758. Nothing inside a retained excerpt is omitted or altered: every retained body is delivered exactly as the article's lines are written, so no omission receipt of any kind accompanies this package. Citations: the retained excerpts carry no citation command at all, so no bibliography entry of the article is transcribed and no reference is redistributed. Reference adaptation: none was needed and none was made; every reference inside the delivered unit resolves inside it, so no unresolved reference or citation is produced. Printed slips kept literal: no printed word, symbol, hypothesis or number of the counted statement or of its proof was altered. Layout and class: the article's own class is \texttt{amsart} with packages including inputenc, amsfonts, amssymb, amsmath, amscd, euscript, verbatim, array, geometry, anysize, fancyhdr, indentfirst, graphicx, color; the wrapper is the lane's pinned \texttt{amsart} editorial wrapper at 10pt on A4, which restates the theorem-like environments the retained text needs and adds the article's own macro definitions that the retained text needs (none), with extra wrapper packages (bm, bbm, dsfont, xspace, cancel, mathtools). The delivered numbering is the unit's own. Assets: no retained span contains a figure environment, a graphics-inclusion command, a PDF-page inclusion, a table or a TikZ picture; no image file is copied into this package, the package declares no asset dependency and no image file is redistributed with it. Licence: version arXiv:2502.10969v3 of this article is distributed under the Creative Commons Attribution 4.0 International licence, as recorded in the arXiv licence metadata for this exact version and shown on the abstract page https://arxiv.org/abs/2502.10969v3. A full rights notice is delivered at licenses/arxiv-2502.10969-CC-BY-4.0.txt. Characters: every retained excerpt is pure ASCII, so the delivered engine, XeLaTeX with its default Latin Modern text font, reports no missing character.
\begin{equation}\label{Z12}
\tilde{N}(\kappa):=n_\kappa-2\gamma_0,
\end{equation}

\begin{Proposition}\label{ko11}
For each $\kappa\in \mathbb{N}$, and each $N\leq \tilde{N}(\kappa)$, there holds \[\bold{K}_N^0(\kappa)\leq 2 \widetilde{\mathbf{K}}_N^1(\kappa,N).\]
\end{Proposition}

\begin{equation}\label{f2x}
\nabla^2(\kappa,\tilde{N},\tilde{N})\leq C_0 q_{\tilde{N}}^{-\frac{\varepsilon}{3}},
\end{equation}

\begin{Lemma}\label{K22}
Given $(v_1,v_2,v_3,v_4)\in (\kappa,r,s)_{II}$, we have
\[K^2(v_1,\ldots,v_4)\lesssim \nabla^2(v_1,\ldots,v_4)+\Theta_4(v_1,\ldots,v_4).\]
\end{Lemma}

\begin{Lemma}\label{keyy1}
Given $\kappa\in\mathbb{N}$, and $0\leq r\leq R\leq \tilde{N}$ where  $\tilde{N}:=\tilde{N}(\kappa)$ is defined by (\ref{Z12}).
\begin{itemize}
	\item [(1)] There holds $\widetilde{\mathbf{K}}_r^1(\kappa,R)\lesssim \varepsilon e^{\frac{\varepsilon}{40}R}q_rq_R^{-1}$.
	\item [(2)] For $j\leq q_r$ and $(v_1,\ldots,v_4)\in (\kappa,R,R)_{II}$, we have
	\begin{align*}
		K^2(j|v_1,\ldots,v_4)-\sum_{i=0}^{j-1}K^2(F^i(v_1,\ldots,v_4))
		\lesssim \varepsilon e^{\frac{\varepsilon}{10}R}q_rq_R^{-1}\sum_{i=0}^{j-1}\frac{\Theta_2(F^i(v_3,v_4))}{\Theta_2(v_3,v_4)}.
	\end{align*}
	\item [(3)] In particular, if $q_R^{1-\frac{\varepsilon}{2}}\sim q_r$, then
	\begin{equation}\label{eeps}
	e^{\frac{\varepsilon}{40}R}q_rq_R^{-1}\leq e^{\frac{\varepsilon}{10}R}q_rq_R^{-1}\sim e^{\frac{\varepsilon}{10}R}q_R^{-\frac{\varepsilon}{2}}<e^{\frac{\varepsilon}{10}R}q_R^{-\frac{\varepsilon}{3}}\leq (\frac{9}{10})^{\frac{\varepsilon R}{2}}.
	\end{equation}
\end{itemize}
\end{Lemma}

\begin{equation}\label{r222}
\sum_{i=0}^{q_r-1}\Theta_2\bigl(F^i(v_1,v_2)\bigr)\leq 2.
\end{equation}

\begin{Proposition}\label{arthh}
Given $\delta\in (0,\frac{\ln 2}{5\ln(1+A_\alpha)})$, for any $m,M\in \mathbb{N}$ with $M\geq m\geq 20$, if $q_M\leq q_m^{1+\delta}$, then $M\leq m+[\frac{m}{2}]\leq \frac{3}{2}m$.
\end{Proposition}

\begin{equation}\label{KES}
\widetilde{\mathbf{K}}_r^1(\kappa,R)\lesssim (\frac{9}{10})^{\frac{\varepsilon R}{2}},
\end{equation}

\begin{Lemma}\label{karrr1} If $q_R^{1-\frac{\varepsilon}{2}}\sim q_r$, then for $r$ large enough,
\begin{equation}\label{karrr}
\Delta(r;R,\kappa)\lesssim (\frac{9}{10})^{\frac{\varepsilon r}{2}}.
\end{equation}
Moreover,
\[\Delta(R;R,\kappa)\lesssim (\frac{9}{10})^{\frac{\varepsilon R}{2}},\quad \widetilde{\mathbf{K}}_R^1(\kappa,R)\lesssim (\frac{9}{10})^{\frac{\varepsilon R}{2}}.\]
\end{Lemma}

\begin{proof}
We prove this result by contradiction. We assume (\ref{karrr}) does not hold.  For fixed $R$ and $\kappa$, we write $\Delta(r)$ instead of $\Delta(r;R,\kappa)$ for simplicity.

At first, we claim that

\noindent{\bf Claim 1: }
For $2< r<R$, and $(v_1,v_2)\in (\kappa,R)_{II}$,
\begin{equation}\label{cla1}
\widetilde{\mathbf{K}}_R^1(\kappa,R)\leq 3\Delta(r).
\end{equation}
\noindent{\bf Proof of Claim 1: }
For $r<R$, $0\leq j<q_R$, break the run of length $j$ into blocks of length $q_r$ and a single remainder block. We obtain
\begin{equation}\label{wan}
{K^1(j|v_1,v_2)\leq \Delta(r)\sum_{i=0}^{j-1}\Theta_2(F^i(v_1,v_2))+\widetilde{\mathbf{K}}_r^1(\kappa,R).}
\end{equation}
If (\ref{karrr}) does not hold, it follows from (\ref{KES}) that the growth of \(\widetilde{\mathbf{K}}_r^1(\kappa,R)\) is negligible compared with \(\Delta(r)\). Together with (\ref{r222}) this yields
\[K^1(j|v_1,v_2)\leq 3\Delta(r),\]
which implies (\ref{cla1}).





\noindent{\bf Claim 2: }
\begin{equation}\label{kaka}
\Delta(R)\lesssim \Delta(r)^2+(\frac{9}{10})^{\frac{\varepsilon R}{2}}.
\end{equation}
\noindent{\bf Proof of Claim 2: }
Let $(v_1,v_2),(v'_1,v'_2)\in (\kappa,R)_{II}$ such that $K^1(q_R|v_1,v_2)$ and $K^1(q_R|v'_1,v'_2)$ have opposite signs. Without loss of generality, we assume \[K^1(q_R|v'_1,v'_2)\leq 0,\]otherwise, we exchange the roles of $v'_1$ and $v'_2$. Then for each $w'>0$, we have
\begin{equation}
|K^1(q_R|v_1,v_2)|\leq K^1(q_R|v_1,v_2)-w'K^1(q_R|v'_1,v'_2).
\end{equation}
Let $i_0\in (0,q_R-1)$ be such that $(F^{i_0}(v_1,v_2),v'_1,v'_2)\in (\kappa,R,R)_{II}$. It follows that for all $i\in \mathbb{Z}$,
\[(F^{i_0+i}(v_1,v_2),F^{i}(v'_1,v'_2))\in (\kappa,R,R)_{II}.\]
Denote
\[i':=\left\{\begin{array}{ll}
	\hspace{-0.4em}i-i_0,\quad i_0\leq i<q_R,\\
	\hspace{-0.4em} i-i_0+q_R,\quad 0\leq i<i_0.
\end{array}\right.\]
Then $F^{i_0+i}(v_1,v_2)$ and $F^{i}(v'_1,v'_2)$ are disjoint. We denote
\[Y(j|v_1,v_2):=\frac{\Theta_2(F^j(v_1,v_2))}{\Theta_2(v_1,v_2)},\]
{\[\Upsilon:=\frac{\Theta_2(F^{i'}(v'_1,v'_2))}{\Theta_2(v'_1,v'_2)}\cdot\frac{\Theta_2(F^{i_0}(v_1,v_2))}{\Theta_2(F^{i}(v_1,v_2))}\]}Then we have
\begin{equation}\label{longf}
\Upsilon=\left\{\begin{array}{ll}
	\hspace{-0.4em}\frac{Y(i'|v'_1,v'_2)}{Y(i'|F^{i_0}(v_1,v_2))},\quad i_0\leq i<q_R,\\
	\hspace{-0.4em} \frac{Y(i'|v'_1,v'_2)}{Y(i'|F^{i_0-q_R}(v_1,v_2))\cdot Y(q_R|F^{-q_R}(v_1,v_2))},\quad 0\leq i<i_0.
\end{array}\right.
\end{equation}
By taking the logarithm of the left hand side of (\ref{longf}) and combining with Proposition \ref{ko11}, we obtain
\[\Upsilon\in \left(\exp(-3\widetilde{\mathbf{K}}_R^1(\kappa,R)),\exp(3\widetilde{\mathbf{K}}_R^1(\kappa,R))\right).\]
By Claim 1, we know that
\[\Upsilon\in \left(\exp(-9\Delta(r)),\exp(9\Delta(r))\right).\]
Choosing
\[w'=\exp(-9\Delta(r))\frac{\Theta_2(F^{i_0}(v_1,v_2))}{\Theta_2(v'_1,v'_2)},\]
Then
\begin{equation}\label{midex}
\exp(-18\Delta(r))\leq w'\frac{\Theta_2(F^{i'}(v'_1,v'_2))}{\Theta_2(F^i(v_1,v_2))}\leq 1.
\end{equation}
Break the runs of length $i_0$ and $q_R-i_0$ into blocks of length $q_r$ plus one shorter block. Denote by $b_p$ the beginning of the $p'$th block. Write
\begin{equation}
\begin{split}
	K^1(q_R|v_1,v_2)-w'K^1(q_R|v'_1,v'_2)=\Xi_1+\Xi_2,
\end{split}
\end{equation}
where
\[ \Xi_1:=\sum_p\left(1-w'\frac{\Theta_2(F^{b'_p}(v'_1,v'_2))}{\Theta_2(F^{b_p}(v_1,v_2))}\right)K^1(q_r|F^{b_p}(v_1,v_2)),\quad \Xi_2:=w'\sum_pB_p,\]
\begin{equation}\label{BB}
B_p:=\left\{\begin{array}{ll}
	\hspace{-0.4em}\Theta_2(F^{b'_p}(v'_1,v'_2))K^2(q_r|F^{b_p}(v_1,v_2),F^{b'_p}(v'_1,v'_2)),\quad j\geq q_r,\\
	\hspace{-0.4em} \Theta_2(F^{b'_p}(v'_1,v'_2))K^2(j|F^{b_p}(v_1,v_2),F^{b'_p}(v'_1,v'_2)),\quad j<q_r.
\end{array}\right.
\end{equation}
%\[w_p:=w'\frac{\Theta_2(F^{b'_p}(v'_1,v'_2))}{\Theta_2(F^{b_p}(v_1,v_2))},\]
By (\ref{midex}), we have
\[\Xi_1\leq (1-\exp(-18\Delta(r)))\Delta(r)\sum_p\Theta_2(F^p(v_1,v_2)).\]
By Lemma \ref{keyy1}(2), we write
\begin{equation}\label{lem6i}\Xi_2\leq \Xi_{21}+\Xi_{22},
\end{equation} where
\[\Xi_{21}:= w'\widetilde{\mathbf{K}}_{r}^1(\kappa,R)\sum_{i=0}^{q_R-1}\Theta_2(F^{i}(v'_1,v'_2)),\]
\[\Xi_{22}:= w'\sum_{i=0}^{q_R-1}\Theta_2(F^{i'}(v'_1,v'_2))K^2(F^{i}(v_1,v_2),F^{i'}(v'_1,v'_2)).\]
By Lemma \ref{keyy1}(1) and (\ref{midex}),
\[\Xi_{21}\lesssim w'\varepsilon(\frac{9}{10})^{\frac{\varepsilon R}{2}}\sum_{i=0}^{q_R-1}\Theta_2(F^{i}(v'_1,v'_2))\leq \varepsilon(\frac{9}{10})^{\frac{\varepsilon R}{2}}\sum_{i=0}^{q_R-1}\Theta_2(F^{i}(v_1,v_2)).\]
Note that $e^{\frac{\varepsilon}{10}R}>1$. By Lemma \ref{K22}, Lemma \ref{keyy1}(3),  (\ref{f2x}) and (\ref{midex}),
\[\Xi_{22}\lesssim e^{\frac{\varepsilon}{10}R}q_R^{-\frac{\varepsilon}{3}}w'\sum_{i=0}^{q_R-1}\Theta_2(F^{i'}(v'_1,v'_2))\lesssim (\frac{9}{10})^{\frac{\varepsilon R}{2}}\sum_{i=0}^{q_R-1}\Theta_2(F^{i}(v_1,v_2)).\]
It follows that
\[\Delta(R)-(1-\exp(-18\Delta(r)))\Delta(r)\lesssim (\frac{9}{10})^{\frac{\varepsilon R}{2}},\]
which from the assumption implies $\Delta(R)<\Delta(r)$ and $\Delta(r)\to 0$ as $r\to \infty$ (up to a subsequence). Moreover, we have (\ref{kaka}). This completes the proof of Claim 2.


The estimate (\ref{karrr}) follows from (\ref{kaka}). In fact, we assume (up to a subsequence)
\[\frac{\Delta(r)^2}{(\frac{9}{10})^{\varepsilon r}}\to +\infty.\]
Due to $q_R^{1-\frac{\varepsilon}{2}}\sim q_r$, it follows from Proposition \ref{arthh} that $r<R<\frac{3}{2}r$. Using (\ref{kaka}) repeatedly, we obtain for any $n\in \mathbb{N}$
\[\frac{\Delta(r)^2}{(\frac{9}{10})^{\frac{\varepsilon r}{2}(\frac{3}{4})^n}}\to +\infty,\]
which contradicts the fact that $\Delta(r)\to 0$ as $r\to \infty$. It follows that
\[\Delta(r)\lesssim (\frac{9}{10})^{\frac{\varepsilon r}{2}},\]
which contradicts our assumption at the beginning of the proof.

Using $r<R<\frac{3}{2}r$ again,  we have
\[\Delta(R)\lesssim (\frac{9}{10})^{\varepsilon r}+(\frac{9}{10})^{\frac{\varepsilon R}{2}}\lesssim (\frac{9}{10})^{\frac{\varepsilon R}{2}}.\]
{It follows from (\ref{wan}) that}
\[\widetilde{\mathbf{K}}_R^1(\kappa,R)\lesssim (\frac{9}{10})^{\frac{\varepsilon R}{2}}.\]
This completes the proof of Lemma \ref{karrr1}.
\end{proof}

\end{document}
